% Suggested replacement/addition to the supporting experimental-validation text.
% Keep existing physical diagnostics separately; no new FE requirement is added.
\paragraph{Recovery of the existing experimental comparison.}
Both a and b in the supplied temperature record denote experimental groups. The three additional worksheet comparisons repeat these observations alongside separately labelled Abaqus histories; they do not provide three further independent experiments. The main figure uses the $20$ mm mesh explicitly identified in the archived validation figure, preserving the prior choice rather than selecting a mesh by its recalculated error. All fourteen plotted FE curves across the available meshes and depths are retained in the accompanying data.

The FE polylines were recovered from vector coordinates calibrated to the original time and temperature ticks. As a reconstruction check, all nine plotted experimental a, b and mean curves were recovered independently and compared with the supplied numerical observations; their maximum discrepancy was $0.852\,{}^\circ$C. This checks plotting consistency and is not a strict bound on all FE polyline simplification errors. Experimental times below $3600$ s were retained, and the $3600$ s endpoint was interpolated within the final recorded segment. FE curves already span $0$--$3600$ s; no temporal extrapolation, initial-temperature shift or calibration was applied.

Time-weighted RMS differences against the group mean are $35.47$, $43.27$ and $39.99\,{}^\circ$C at $15$, $30$ and $40$ mm, respectively. Differences against a and b separately are retained, as are the original finite-element plot and observation-source identities. Exact group-to-specimen assignments, thermocouple positions and the complete original solver histories remain unavailable. The recovered comparison supplies the existing experimental support for the thermal modelling route and is distinct from the subsequent numerical sensitivity calculations.
